\scanpage{14}

\uline{Exemple}

\uline{Corollaire} Soit $G$ un préschéma en groupes de \deleted{type}\inserted{prés} \unsure{fini|finie} sur $S$, \deleted{loc. noeth.} $P$ un \deleted{fibré principal} préschéma \unsure{de|\deleted{de}} prés. finie sur $S$ dans lequel $G$ opère, de façon que $P$ soit form. pr. hom. sous $G$. \unsure{Dans|Sous} ces conditions, il existe un mono de préschémas $S'\hookrightarrow S$, tel que $P\longrightarrow S$ se factorise par $P\unsure{\text{---}|\to}S'$, et que \deleted{\unsure{Reste} $P'$} $P$ devienne \deleted{formellement} principal hom sous $G_{S'}$ i.e $P\to S'$ est \deleted{plat} \uline{fid.} \uline{plat}. [De \deleted{\unsure{tt}} plus, $S'$ est \unsure{uniqt|unique,|uniquement|unique} déterminé à $S$-isom près, et est de prés. finie sur $S$].

\uline{Corollaire} Soit $u:G\unsure{\to|\text{---}}H$ un hom de $S$-préschémas en groupes, \deleted{de présentation} $u$ de présentation finie. Supposons $\mathrm{Ker}\,u=\deleted{\unsure{K}}N$ [qui est de prés finie sur $S$] \uuline{plat} sur $S$. \uline{Alors} \deleted{\unsure{$u$ se factorise|le foncteur}} $G/N$ \uline{existe}.
% A vertical pen line marks the second corollary.
% archive sheet number: 13
